\begin{mistake}{2026-04-12}{29}

\begin{problemcontent}
已知 \( f_1(x), f_2(x) \) 是二阶线性齐次微分方程 \( y''+py'+qy=0 \) 的两个特解，求使 \( C_1f_1(x)+C_2f_2(x) \) 成为该方程通解的条件。
\end{problemcontent}

\begin{solutioncontent}
二阶线性齐次微分方程的解空间维数为 2，因此要使
\[
y=C_1f_1(x)+C_2f_2(x)
\]
成为通解，关键在于 \(f_1,f_2\) 必须构成一组基本解，即二者线性无关。

判定线性无关可用朗斯基行列式：
\[
W(f_1,f_2)(x)=
\begin{vmatrix}
f_1(x) & f_2(x) \\
f_1'(x) & f_2'(x)
\end{vmatrix}
=f_1(x)f_2'(x)-f_2(x)f_1'(x)
\]
因此所求条件是
\[
W(f_1,f_2)(x)\neq 0
\]
在所讨论区间内至少对某一点成立；等价地说，\(f_1,f_2\) 在线性意义下无关。

于是，只有当 \(f_1,f_2\) 线性无关时，\(C_1f_1(x)+C_2f_2(x)\) 才能表示该方程的全部解，也即成为通解。
\end{solutioncontent}

\mistakeknowledge{
\begin{itemize}
  \item \textbf{核心公式}：朗斯基行列式 \(W(f_1,f_2)=f_1f_2'-f_2f_1'\)；二阶线性齐次方程的通解由两组线性无关特解线性组合得到。
  \item \textbf{易错点}：不要把“是两个特解”误认为“必能组成通解”；还必须检验线性无关。二阶行列式展开顺序也常写反。
  \item \textbf{关联知识}：对 \(n\) 阶线性齐次方程，需要 \(n\) 个线性无关解才能组成通解，朗斯基行列式是标准判定工具。
\end{itemize}}

\end{mistake}
